
\section{Status of the Speculative Program}\label{sec:part2_preamble}

Part~II takes structures built in Part~I (agent frames, transport, the consensus free energy, the lagged predictive prior supplied by a parent) and asks what they would mean if read physically: as time, as geometry, as observer-relative reality. Its claims have status S (structural correspondence or interpretation) or C (conjectured or open) in the scheme of Table~\ref{tab:claim_status}, and each rests on one of the numbered postulates P1 to P10 below. None of it makes a quantitative prediction, and in its central domains (the Lorentzian signature, gravity, consciousness) the material is interpretation, not falsifiable physics. The mathematics that supports the readings has its own status and does not depend on them: the clock-dissipation identity of Eq.~\eqref{eq:clock_dissipation} (P, for an autonomous single-rate flow), the binary registration bound of Proposition~\ref{prop:binary_registration} (D), the bit-production ratio of Eq.~\eqref{eq:bit_production_ratio} (D), the covariant pullback of Eq.~\eqref{eq:induced_metric_full} and the passive gauge invariance of its covariant Fisher block (D for the pullback, P for the invariance, which is proved in the companion Supplement~B~\cite[Ch.~B2]{DennisGaugeVFESuppB2026}), the signature constructions of Section~\ref{sec:signature_resolution} (P under their postulates), and the regulator obstruction of Section~\ref{sec:consensus_metric} (P), which concerns only the frame-twist block. An existence demonstration of structural compatibility is never presented as a derivation. Whether the participatory reading of the formalism admits a decidable test is recorded in Table~\ref{tab:claim_status}; the physical readings of Part~II do not.

The companion theory declines, without refuting, each physical identification made below~\cite[Chs.~1 and~14]{DennisGaugeVFE2026} (Section~\ref{sec:companion_theory}); Part~II adopts these identifications as the postulates P1 to P10. The extensions stay in this paper because each bears on a question Wheeler's participatory program raises: how duration, geometric signature, observer-relativity, and dimensional structure relate to information-geometric data. Read constructively, Part~II fixes the program's next constructions in the order C1 through C6 of Section~\ref{sec:open_problems}.

The postulate register follows. Every entry is a postulate, not a consequence of free-energy minimization, and the sections below cite entries by number. The detailed assumptions of the algebraic signature route are numbered P4.1 to P4.6 as sub-items of entry P4, so no second numbering scheme arises.
\begin{description}
\item[P1 (clock).] An agent's internal time is the Fisher-Rao arc length $\vartheta_i$ along its belief history, whose order and orientation are supplied by the shared flow parameter $t$; the bit count of a continuous update is a fixed-granularity diagnostic, and Wheeler's bit is represented separately by binary registrations (Section~\ref{sec:time_info}).
\item[P2 (bundle metric).] The associated bundle carries the fiber-respecting tensor with vertical block the covariant Fisher-Rao pullback and horizontal block $\kappa\operatorname{tr}(A_\mu A_\nu)$ with a declared coupling $\kappa$ (Section~\ref{sec:it_from_bit}).
\item[P3 (tier identification).] Perceived ambient spatial structure is the structural-tier pullback $G^{(s)}$ (Section~\ref{sec:three_tiers}).
\item[P4 (signature, algebraic route).] On a declared reference section, a commuting, mutually trace-orthogonal generator set with one temporal direction of negative trace-square and three spatial directions of positive trace-square (all-real route); on the complex route, instead, a single generator of nonzero trace-square together with the bilinear pairing $\operatorname{tr}(AB)$, the imaginary temporal coefficient, and the real-part projection (Postulates~P4.1 to P4.6 of Section~\ref{sec:signature_resolution}).
\item[P5 (signature, causal route).] A coarse-grained shared update parameter $t$ used as a base coordinate, a positive-definite spatial information metric $h_{ab}$, an influence region that is an $h$-ball, and a finite maximal epistemic propagation speed $c_\mathcal{I}$ (Section~\ref{sec:causal_cone_route}).
\item[P6 (sector hierarchy).] The eigenvalue hierarchy $\lambda_{\text{obs}} \gg \lambda_{\text{subthresh}} \gg \lambda_{\text{internal}}$ with constant sector ranks, and $|\mathcal{D}_{\text{obs}}| = 4$ for human agents, which requires $\dim\mathcal{C} \ge 4$ (Section~\ref{sec:observable_sectors}).
\item[P7 (dimensional monism).] Nats are dimensionally fundamental; physical dimensions are consensus labels (Section~\ref{sec:phenomenological_interpretation}).
\item[P8 (structural curvature).] Gravitational phenomenology, if any, is curvature of the structural pullback $G^{(s)}$; the equivalence principle requires an undelivered link between $\Sigma_p^{-1}$ and $\Sigma_s^{-1}$ (Section~\ref{sec:gravity_pullback}).
\item[P9 (phenomenal correlate).] Phenomenal character corresponds to the geometric structure of the induced metric, with the frame-to-quale assignment conjectured fixed by the cross-scale composition (Section~\ref{sec:consciousness}).
\item[P10 (shareability).] Scientific content must be statable frame-independently; this at most motivates general covariance (Section~\ref{sec:gauge_invariance_consensus}).
\end{description}

\section{Time as Registered Change: The Information Clock}\label{sec:time_info}

This section develops W1 of Table~\ref{tab:wheeler_correspondence}, registered change as a clock. The dynamics of Part~I are written in a flow parameter $t$ shared by all agents, the continuous limit of the update index, and $t$ supplies the order and the orientation of every belief history. No physical time is presupposed, but an ordering is: the construction measures each agent's progress along a common ordering and does not recover the ordering from correlations, which is the step that Page--Wootters and thermal-time proposals take. The clock $\vartheta_i$ of agent $i$ is the Fisher-Rao arc length of its belief history, an additive length that is invariant under refinement and reparameterization and has a unit-speed form~\cite[Ch.~7]{DennisGaugeVFE2026}. Postulate P1 reads it as the agent's internal time, an identification the companion theory declines~\cite[Ch.~1]{DennisGaugeVFE2026}.

A belief update carries the information gain $\Delta I_i = D_{\mathrm{KL}}(q_i^{\text{new}} \| q_i^{\text{old}})$, in nats (division by $\ln 2$ converts to bits), and a step of Fisher length $\delta\vartheta$ carries $\delta\vartheta^2/(2\ln 2)$ bits to leading order by Eq.~\eqref{eq:kl_arclength} below. This count measures belief change in bits; it does not model Wheeler's elementary yes-no registration, which Section~\ref{sec:binary_registration} treats separately. The unit of change is set by statistical distinguishability: two beliefs a parameter step $\delta\theta$ apart are distinguishable at the level of one standard error from $N$ independent observations when $N g_{\mathcal{B}}(\delta\theta, \delta\theta) \gtrsim 1$, with $g_{\mathcal{B}}(\delta\theta,\delta\theta)$ the squared Fisher-Rao length~\cite{Wootters1981, BraunsteinCaves1994, AmariNagaoka2000}. This is Wootters' statistical distance, which Wheeler cited as a candidate link between information and geometry~\cite{Wheeler1990}, and it is the information-geometric reading of Bateson's ``a difference which makes a difference''~\cite{Bateson1972}. For a Gaussian belief and a pure mean step at fixed covariance,
\begin{equation}
D_{\mathrm{KL}}(q^{\text{new}} \| q^{\text{old}}) = \tfrac{1}{2}\delta\mu^\top \Lambda \delta\mu = \tfrac{1}{2}ds^2,
\label{eq:kl_arclength}
\end{equation}
exactly, with $\Lambda = \Sigma^{-1}$ the precision and $ds^2$ the squared length of the straight constant-covariance mean path; the Fisher-Rao distance of the full Gaussian family between the same endpoints can be shorter~\cite[Ch.~7]{DennisGaugeVFE2026}. Because the increment is quadratic in the step, the bit count is not additive under refinement: a path of unit Fisher length taken in one step carries $0.7213$ bits, and the same path taken in $100$ equal substeps carries $0.0072$ bits in total, while its arc length stays $1$. The bit count is therefore defined only relative to a fixed update granularity, which is why P1 takes the arc length as primary.

Different agents accumulate different arc lengths. No time function common to all agents is built into the construction. Along a history of the population with block-diagonal metric the squared speeds add, so the collective clock $\vartheta$ obeys $\max_i \vartheta_i \le (\sum_i \vartheta_i^2)^{1/2} \le \vartheta \le \sum_i \vartheta_i$; a clock that is a function of the population state and orthogonal to the descent direction exists only under an exactness condition that fails generically, although unit-speed clocks built from a global cross-section can exist; and the clock of a gauge-invariant free energy is itself gauge invariant, since every natural-gradient history is then orthogonal to the gauge orbits~\cite[Ch.~7]{DennisGaugeVFE2026}. The closest precedent is Caticha's entropic time, in which time is introduced as a device to keep track of change in an information-updating dynamics~\cite{Caticha2011entropic}; what the present construction adds is a clock per agent, read in that agent's frame, and the meta-agent clock of Section~\ref{sec:scale_dependent_time}. The clock also differs from proposals that read time off an entangled clock subsystem~\cite{PageWootters1983}, off the modular flow of the system's own global statistical state~\cite{ConnesRovelli1994}, off computational steps~\cite{Lloyd2002}, or off horizon thermodynamics and entanglement structure~\cite{Jacobson1995, VanRaamsdonk2010}: here it is the length of each agent's own belief history~\cite{Wheeler1990, Rovelli2004}.

\subsection{Fisher Arc Length Along Belief Trajectories}
\label{sec:fisher_arc_length}

An agent following a belief trajectory $q(\lambda)$, with $\lambda$ any orientation-preserving parameter, accumulates the Fisher-Rao arc length $\Delta\vartheta = \int \sqrt{g_{\mathcal{B}}(q'(\lambda), q'(\lambda))} d\lambda$, a reparameterization-invariant Riemannian length, and the clock $\vartheta$ is the arc-length parameter it defines. This clock is not relativistic proper time. Proper time is the Lorentzian length of a timelike worldline and is largest for an unaccelerated clock at rest, whereas the Fisher clock accumulates only while the belief changes and is stationary for an unchanging belief: it is a change-based duration. Any connection to physical time dilation would require the indefinite-metric extension of Section~\ref{sec:signature_resolution}, and none is established here.

The clock and the dynamics are linked by an exact identity. Let the belief sector follow an autonomous, single-rate Fisher natural-gradient flow $\dot{x} = -\eta G^{-1}\nabla\mathcal{F}$, with $G$ the Fisher metric and $\eta$ one scalar rate. Then $G(\dot{x}, \dot{x}) = \eta^2 \nabla\mathcal{F}^\top G^{-1}\nabla\mathcal{F}$ and $d\mathcal{F}/dt = \nabla\mathcal{F}^\top \dot{x} = -\eta \nabla\mathcal{F}^\top G^{-1}\nabla\mathcal{F}$, hence
\begin{equation}
\left(\frac{d\vartheta}{dt}\right)^2 = \eta\left(-\frac{d\mathcal{F}}{dt}\right),
\label{eq:clock_dissipation}
\end{equation}
so along such a flow the information clock ticks exactly while free energy is dissipated. The identity is the natural-gradient counterpart of the thermodynamic-length relations between Fisher speed and dissipation~\cite{SalamonBerry1983, Crooks2007, Ito2018}, and the companion theory states it in rate form together with its reparameterization law~\cite[Ch.~7]{DennisGaugeVFE2026}. For one discrete step $x_{k+1} = x_k - \eta G^{-1}\nabla\mathcal{F}$ with dimensionless step size $\eta$, the information gain of Eq.~\eqref{eq:kl_arclength}, in nats, is $\Delta I = (\eta/2)|\Delta\mathcal{F}|$ to leading order in $\eta$.

Three qualifications fix the scope of Eq.~\eqref{eq:clock_dissipation}. First, when the free energy depends explicitly on time through the data, as for an agent receiving observations $o(t)$, $d\mathcal{F}/dt = \nabla\mathcal{F}^\top\dot{x} + \partial_t\mathcal{F}$ and the identity becomes
\begin{equation}
\left(\frac{d\vartheta}{dt}\right)^2 = \eta\left(\partial_t\mathcal{F} - \frac{d\mathcal{F}}{dt}\right),
\label{eq:clock_driven}
\end{equation}
so the clock can tick while $\mathcal{F}$ rises. Second, the belief flow of Eq.~\eqref{eq:gauge_natural_gradient} carries separate rates $\eta_\mu$ and $\eta_\Sigma$; for such a rate-separated flow the identity holds block by block and holds with unit rate in the rate-weighted metric $\bigoplus_b G_b/\eta_b$, while the unweighted Fisher clock obeys only the bound $|d\mathcal{F}/d\vartheta| \le \|\operatorname{grad}^{G}\mathcal{F}\|_{G}$~\cite[Chs.~7 and~8]{DennisGaugeVFE2026}. The frame retraction of Eq.~\eqref{eq:gauge_group_retraction} is preconditioned by a metric that is not Fisher, so frame motion contributes a length but not an information clock. Third, for a real-analytic free energy an autonomous flow that converges has finite total clock reading, whereas a smooth free energy admits a converging flow of infinite duration; the companion theory proves the first statement and exhibits the counterexample~\cite[Ch.~7]{DennisGaugeVFE2026}. Sustained clock advance therefore requires sustained driving by observations, and Eq.~\eqref{eq:clock_driven}, not Eq.~\eqref{eq:clock_dissipation}, governs an observing agent. On the P1 reading, an agent's internal time along an autonomous single-rate flow is accumulated free-energy dissipation measured in the Fisher metric, and along a driven flow it is that dissipation plus the explicit change of $\mathcal{F}$ produced by the data.

\subsection{Binary Registrations and Wheeler's Bit}
\label{sec:binary_registration}

Wheeler's bit is an event, a registered answer to a yes-no question elicited by a device~\cite{Wheeler1990}, not a unit of measure on a continuous update. The framework can represent that event directly by an observation channel with binary outcomes, and the expected belief change it produces is then bounded by one bit and additive over registrations.

\begin{proposition}[Information budget of binary registrations]\label{prop:binary_registration}
Let $X$ be a latent variable with prior law $p(x)$ and let $O \in \{0,1\}$ be a yes-no observation with likelihood $P(O=1 \mid x)$. The expected belief update equals the mutual information and is at most one bit,
\begin{equation}
\mathbb{E}_{O}\left[D_{\mathrm{KL}}\left(p(x \mid O) \| p(x)\right)\right] = I(X;O) \le H(O) \le \ln 2 .
\label{eq:registration_bound}
\end{equation}
For registrations $O_1, \ldots, O_n$, the expected updates add by the chain rule, $I(X; O_1, \ldots, O_n) = \sum_{k=1}^n I(X; O_k \mid O_{<k})$, where the $k$-th term is the expected update produced by $O_k$ from the belief $p(x \mid O_{<k})$. If the registrations are conditionally independent given $X$, each term is at most $I(X;O_k) \le \ln 2$, so $n$ registrations change the belief by at most $n$ bits in expectation.
\end{proposition}
\begin{proof}
By definition,
\begin{align*}
\mathbb{E}_O D_{\mathrm{KL}}(p(x\mid O)\|p(x)) &= \sum_o P(o)\int p(x\mid o)\ln\frac{p(x\mid o)}{p(x)}dx \\
&= \int\sum_o p(x,o)\ln\frac{p(x,o)}{p(x)P(o)}dx = I(X;O).
\end{align*}
Since $O$ is discrete, $H(O \mid X) \ge 0$, so $I(X;O) = H(O) - H(O\mid X) \le H(O) \le \ln 2$. The chain rule for mutual information~\cite[Thm.~2.5.2]{CoverThomas2006} gives the sum, and applying the first part to the conditional laws given $O_{<k}$ identifies each term as an expected update. Under conditional independence, $I(O_k; X, O_{<k}) = I(O_k; X)$, and expanding the left side by the chain rule gives $I(X;O_k \mid O_{<k}) = I(X;O_k) - I(O_k;O_{<k}) \le I(X;O_k)$.
\end{proof}

The bound holds in expectation only: a single surprising outcome can move a belief by more than one bit. For a standard normal prior and a sharp logistic channel whose positive answer has probability $1.4\times10^{-3}$, the positive outcome moves the belief by $9.36$ bits, while the expected update is $0.0145$ bits; for symmetric channels of low and high sharpness the expected update is $0.135$ and $0.962$ bits. The registration count, and the information budget it bounds, are additive, whereas the Fisher clock of Section~\ref{sec:fisher_arc_length} is a path length whose continuous-update bit count depends on the granularity of the update. On this construction a continuous Gaussian belief is an idealization of accumulated registrations, which is Wheeler's order of derivation; that reading is status S. The channel models the inferential half of Wheeler's observer-participant only. No agent chooses which question to pose, the registration is not irreversible in any physical sense, and a classical fiber has no complementarity, so ``no question, no answer'' and the elementary quantum phenomenon remain outside the model.

\subsection{The Scale-Dependent Meta-Agent Clock}
\label{sec:scale_dependent_time}

A meta-agent that aggregates many constituents does not advance its clock at the constituent rate, because a difference registers on the meta-clock only to the extent that it is shared across the cluster whose gauge-covariant barycenter the meta-belief is. Decompose the constituent mean updates into common and idiosyncratic parts, $\delta\mu_i = \delta c + \epsilon_i$, with $\mathbb{E}[\epsilon_i] = 0$, the $\epsilon_i$ uncorrelated with one another and with $\delta c$, and a common second moment $\mathbb{E}\|\epsilon_i\|^2_\Lambda = \mathbb{E}\|\epsilon\|^2_\Lambda$, and define the coherent fraction $f_{\mathrm{coh}} := \mathbb{E}\|\delta c\|^2_\Lambda / (\mathbb{E}\|\delta c\|^2_\Lambda + \mathbb{E}\|\epsilon\|^2_\Lambda)$ in the precision norm $\|v\|^2_\Lambda = v^\top \Lambda v$.
\begin{proposition}[Meta-clock slowdown]\label{prop:meta_clock_slowdown}
In the uniform-weight mean-sector special case of the coarse-graining construction (all coherence weights equal in Eq.~\eqref{eq:meta_agent_mu_impl}, with the averaging covariance convention of Eq.~\eqref{eq:meta_agent_sigma_impl}), the meta-to-constituent ratio of expected per-step bit production is
\begin{equation}
Z_I := \frac{\mathbb{E}[\Delta I_I]}{\mathbb{E}[\Delta I_i]} = f_{\mathrm{coh}} + \frac{1 - f_{\mathrm{coh}}}{N} = f_{\mathrm{coh}} + O(1/N),
\label{eq:bit_production_ratio}
\end{equation}
with $\Delta I$ the information gain of Eq.~\eqref{eq:kl_arclength}.
\end{proposition}
\begin{proof}
Under the averaging convention the meta-agent's precision is $\Lambda$ and its mean update is $\delta c + N^{-1}\sum_i \epsilon_i$, so by Eq.~\eqref{eq:kl_arclength} and the stated correlations $2\mathbb{E}[\Delta I_I] = \mathbb{E}\|\delta c\|^2_\Lambda + N^{-1}\mathbb{E}\|\epsilon\|^2_\Lambda$, while $2\mathbb{E}[\Delta I_i] = \mathbb{E}\|\delta c\|^2_\Lambda + \mathbb{E}\|\epsilon\|^2_\Lambda$.
\end{proof}
The ratio runs from $Z_I \to 1$ for fully coherent activity to $Z_I \to 1/N$ for fully incoherent activity, so registering one meta-tick on the bit clock requires of order $N$ bits of constituent activity. The bit is a fixed reference unit held constant across scales; the slowdown resides in the rate at which bits are produced, and the $1/N$ suppression is the law of large numbers acting on the barycenter, not a dynamical critical exponent. The proposition is stated in the relative-entropy currency. Its square root $\sqrt{Z_I}$ is the ratio of root-mean-square Fisher speeds, not the ratio of expected arc lengths, because the expected ratio of arc lengths is not the square root of the expected ratio of relative entropies; the arc-length ratio on the primary clock of P1 is an open problem of the companion theory~\cite[Ch.~7]{DennisGaugeVFE2026}, of which Proposition~\ref{prop:meta_clock_slowdown} answers the relative-entropy version.

A comparison that does not depend on the covariance convention is available. In this special case the meta-agent law is the image of the constituents' joint law under a fixed Markov kernel for either convention: the averaging convention is the image under $y = N^{-1}\sum_i x_i + \xi$ with independent Gaussian noise $\xi$ of covariance $\Sigma - \Sigma/N$, and the precision-pooling convention $\Lambda_I = N\Lambda$, which the meta-agent state construction declines to adopt (Section~\ref{sec:meta_agent_state_construction}), is the deterministic averaging image. By the same-path record-clock contraction of the companion theory, a law that is the image of a history under one parameter-independent Markov kernel accumulates no more Fisher duration than that history, with defect equal to the expected conditional variance of the score~\cite[Ch.~7]{DennisGaugeVFE2026}. The meta-agent is therefore never faster than the collective of its constituents under either convention. Against the joint history of the $N$ constituents, the precision-pooled parent has expected squared-speed ratio $f_{\mathrm{coh}} + (1 - f_{\mathrm{coh}})/N$, and the averaging parent has that ratio divided by $N$. Only the comparison with a single constituent depends on the convention: the averaging convention gives the slowdown $Z_I$ of Eq.~\eqref{eq:bit_production_ratio}, and precision pooling gives $N f_{\mathrm{coh}} + (1 - f_{\mathrm{coh}})$, a speedup relative to one constituent. Each of these ratios follows from the second-moment computation in the proof of Proposition~\ref{prop:meta_clock_slowdown}. The ratio itself follows from the consensus coupling already present in the free energy: to leading order the belief-coupling term reduces to a graph-Laplacian quadratic form $\mu^\top (L \otimes \Lambda)\mu$ in which the barycenter is the zero mode that takes a $1/N$ share of incoherent injected power, while the algebraic connectivity $\lambda_2$ of the coupling graph~\cite{Fiedler1973} sets the rate $2\eta_\mu\lambda_2$ at which constituent disagreement relaxes and the detector of Section~\ref{sec:meta_agent_threshold} triggers.

\section{Observer-Dependent Geometry: The Pullback Construction}\label{sec:it_from_bit}

\subsection{The Pullback Mechanism: From Belief Sections to Metric Data}
\label{sec:pullback}

Wheeler proposed, as a working hypothesis, that every physical quantity, ``even the spacetime continuum itself'', derives its significance from binary yes-or-no indications elicited by observing devices~\cite{Wheeler1990}. This section develops W2 of Table~\ref{tab:wheeler_correspondence}: a classical information-geometric model in which metric data on a base that carries no physical metric are derived from agents' belief sections, so that geometry is observer-relative. The pullback tensor has status D and the gauge invariance of its covariant Fisher block status P; their identification with physical geometry is status S and rests on postulates P2 and P3. The geometry is derived from continuous belief sections, not from discrete registrations; Section~\ref{sec:binary_registration} gives the registration channel.

The base manifold $\mathcal{C}$ is a presupposed smooth index manifold carrying no physical metric. Its topology, dimension, and smooth structure are inputs, and only metric data are derived, so the construction models Wheeler's thesis for the metric alone and not his rejection of a presupposed continuum, the ``no continuum'' among his four no's~\cite{Wheeler1990}. The literature indicates a route that would address the gap: for a distinguishing spacetime already known to be Lorentzian, the causal order fixes topology, smooth structure, and conformal metric~\cite{HawkingKingMcCarthy1976, Malament1977}, and causal set theory recovers the volume element by counting~\cite{BombelliLeeMeyerSorkin1987}. Building the base from agent data along such lines is work on the base itself and is deferred to the companion Supplement~B~\cite{DennisGaugeVFESuppB2026}.

Agent $i$ carries its belief as a section $q_i$ over $\mathcal{C}$, expressed in its frame $U_i(c)$, with Regime~I connection $A^{(i)}_\mu = U_i^{-1}\partial_\mu U_i$ and field strength $F_{\mu\nu} = 0$ identically. The covariant first jet of the section is
\begin{equation}
D_\mu q_i = \partial_\mu q_i - \rho_*\left(\partial_\mu U_i U_i^{-1}\right) q_i ,
\label{eq:covariant_jet}
\end{equation}
which for a Gaussian section reads $D_\mu(\mu_i, \Sigma_i) = (\partial_\mu\mu_i - X_\mu\mu_i, \partial_\mu\Sigma_i - X_\mu\Sigma_i - \Sigma_i X_\mu^\top)$ with $X_\mu = \partial_\mu U_i U_i^{-1}$. Equivalently, $\rho(U_i^{-1})D_\mu q_i = \partial_\mu \hat{q}_i$ for the frame-stripped section $\hat{q}_i = \rho(U_i^{-1})q_i$, and because $\rho(U_i^{-1})$ is a Fisher isometry the pullback below is the ordinary pullback of $\hat{q}_i$. Writing $\nabla^{(i)}_\mu \log q_i$ for the score of $D_\mu q_i$, the covariant Fisher pullback is
\begin{equation}
\boxed{G^{(q)}_{i,\mu\nu}(c) = \mathbb{E}_{q_i(c)}\big[(\nabla^{(i)}_\mu \log q_i)(\nabla^{(i)}_\nu \log q_i)\big],}
\label{eq:induced_metric_full}
\end{equation}
a positive semi-definite tensor whose rank is the rank of the jet~\cite[Ch.~B2]{DennisGaugeVFESuppB2026}. For constant frames $D_\mu = \partial_\mu$ and Eq.~\eqref{eq:induced_metric_full} is the ordinary pullback of the Fisher-Rao metric. Postulate P2 adds a horizontal frame-twist block with a declared coupling $\kappa$,
\begin{equation}
\widetilde{G}^{(q)}_{i,\mu\nu}(c) = \kappa\operatorname{tr}\big(A^{(i)}_\mu A^{(i)}_\nu\big) + G^{(q)}_{i,\mu\nu}(c).
\label{eq:combined_tensor}
\end{equation}
The coupling is needed because the two blocks are unrelated in normalization: the trace block scales with the normalization of the Lie-algebra generators, the Fisher block does not, and without $\kappa$ the signature of the sum would depend on an arbitrary generator normalization (Section~\ref{sec:signature_resolution}). The trace form is Ad-invariant but indefinite on real $\mathfrak{gl}(K_q)$: for $J=\left(\begin{smallmatrix}0&-1\\1&0\end{smallmatrix}\right)$, $\operatorname{tr}(J^2)=-2$, whereas for $D=\operatorname{diag}(1,0)$, $\operatorname{tr}(D^2)=1$. The horizontal block is therefore an indefinite trace tensor, not a Riemannian metric, and the compact alternative $-\operatorname{tr}(AB)$ is positive only on a skew-symmetric compact subalgebra. The construction transfers to the prior, generative-model, and hyper-prior sections, yielding $G_i^{(q)}, G_i^{(p)}, G_i^{(s)}, G_i^{(r)}$.

The two blocks behave differently under the redundancies of the functional. Under a passive local re-expression of agent frames, $U_i \mapsto g_i(c)U_i$ with $q_i \mapsto \rho(g_i(c))q_i$ and $\Omega_{ij} \mapsto g_i\Omega_{ij}g_j^{-1}$, the frame-stripped section $\hat{q}_i$ is unchanged, so $G^{(q)}_i$ is invariant with no orbit average; this is the passive gauge covariance of covariant pullbacks proved in the companion Supplement~B~\cite[Ch.~B2]{DennisGaugeVFESuppB2026}, where the pullback is also shown to be relative to the chosen connection, here the connection defined by the declared reference section common to all agents. The ordinary fixed-frame pullback of $q_i$ is not invariant under this re-expression, since it acquires frame-gradient terms, and the frame-twist block is invariant under neither local action: the change of reference $U_i(c) \mapsto U_i(c)g(c)$ leaves every $\Omega_{ij}$, and hence the free energy, unchanged, yet with $g = U_i^{-1}$ it sets $\operatorname{tr}(A^{(i)}A^{(i)})$ to zero for any chosen agent. At a fixed reference section, then, the differences $G^{(q)}_i \neq G^{(q)}_j$ between agents are differences between their beliefs expressed in a common reference, not differences of frame choice, and the regulator obstruction of Section~\ref{sec:consensus_metric} concerns only the frame-twist block. The covariant construction uses base derivatives of sections and frames and so belongs to the base-transport material of Supplement~B; within the single-base-point scope of the companion theory the base tangent space is zero and every pullback tensor vanishes.

Both tensors are dimensionless: the Fisher pullback is invariant under a change of fiber units and scales quadratically under a rescaling of base coordinates, so $G^{(q)}_{i,\mu\nu}dc^\mu dc^\nu$ is a pure number of nats in every chart, and the frame-twist block is dimensionless because the frame fields are Lie-algebra coordinates. A physical line element would require a length scale $L_0$, with $ds^2_{\text{phys}} = L_0^2 G$, and a Lorentzian reading would further require a constant conversion between timelike and spacelike normalizations; the algebraic route of Section~\ref{sec:signature_resolution} supplies neither, since its cone slope is field-dependent, and the causal route postulates the speed.

The pulled-back tensors enter no term of the free energy; the integration measure on $\mathcal{C}$ is the background chart density. The construction therefore supplies the arc of Wheeler's loop from information to geometry but not the arc from geometry back to the participants. Closing the loop would require the induced geometry to enter the dynamics, for example as the integration measure or as the coupling neighborhood, which is base-transport work deferred to Supplement~B. Information-geometric models of physical space and of spacetime were proposed by Caticha~\cite{Caticha2005, Caticha2019spacetime}, and Fisher-metric constructions of an Einstein tensor exist~\cite{Matsueda2013}; the lineage runs back to Wootters' statistical distance~\cite{Wootters1981}. The contribution here is the agent-indexed family of pullbacks and its gauge behavior, not the idea of Fisher geometry as spatial geometry.

\subsection{Three Tiers of Induced Geometry: Epistemic, Expectational, Structural}
\label{sec:three_tiers}

The four induced tensors carry different interpretations, organized by the timescale on which their underlying section evolves. The belief pullback $G_i^{(q)}$ is epistemic geometry, an instantaneous uncertainty map fluctuating on the perceptual timescale. The prior pullback $G_i^{(p)}$ is expectational geometry: its section is the prior supplied by the parent through the lagged predictive prior of Section~\ref{sec:cross_scale_shadows}, which for a transported-copy transition is $p_i = \Omega_{iI}[q_I^{(\ell+1)}]$ from the previous episode, so in its covariant form it inherits the profile of the parent's belief and evolves more slowly than the epistemic tier while remaining a state-fiber object. The generative-model pullback $G_i^{(s)}$, with its hyper-prior partner $G_i^{(r)}$, is structural geometry on the model fiber, quasi-static on the learning timescale. Postulate P3 identifies what an agent experiences as ambient spatial structure with the structural tier $G_i^{(s)}$. The identification is chosen by timescale matching: perceived ambient space is invariant across changing percepts on perceptual timescales, and among the four tiers only the model-fiber pair is quasi-static on that timescale, with $G_i^{(s)}$ the member carrying learned structure. The reading echoes Kant's space as a form of intuition, Wheeler's observer-participatory universe, and Hoffman's interface theory of perception~\cite{Hoffman2019}; it places the interface in the slow generative model, whereas predictive-processing accounts usually identify perceptual content with the posterior~\cite{Clark2016, Seth2021}. Agents with different generative models therefore carry different structural geometries on the same base, and within this interpretive frame there is no agent-independent metric on $\mathcal{C}$, only the family of agent-relative induced tensors. A control in which the slow channel is frozen would make $G_i^{(s)}$ and $G_i^{(r)}$ constant and leave the dynamical content on the fast channel; no such experiment is reported here.

\subsection{Example: Gaussian Agents on \texorpdfstring{$\mathbb{R}^2$}{R2}}

For agents with constant frames maintaining Gaussian beliefs $q_i(c) = \mathcal{N}(\mu_i(c), \Sigma_i(c))$ over a flat base $\mathcal{C} = \mathbb{R}^2$, the induced tensor is
\begin{equation}
G^{(q)}_{i,\mu\nu}(c) = (\partial_\mu \mu_i)^\top \Sigma_i^{-1} (\partial_\nu \mu_i) + \tfrac{1}{2}\operatorname{tr}\left(\Sigma_i^{-1}(\partial_\mu \Sigma_i)\Sigma_i^{-1}(\partial_\nu \Sigma_i)\right),
\end{equation}
whose first term captures mean-gradient and second term covariance-gradient contributions; with nonconstant frames the derivatives become the covariant jets of Eq.~\eqref{eq:covariant_jet}. In the isotropic case $\Sigma_i(c) = \sigma(c)^2 I$ it reduces to
\begin{equation}
G^{(q)}_{i,\mu\nu}(c) = \sigma^{-2}\left[(\partial_\mu\mu_i)\cdot(\partial_\nu\mu_i) + 2K_q (\partial_\mu\sigma)(\partial_\nu\sigma)\right],
\end{equation}
positive definite only where the jet $(d\mu_i, d\sigma)$ has rank two. The factor $\sigma^{-2}$ acts as a conformal factor on the mean term, so where $\sigma$ varies, confident regions appear larger and uncertain regions smaller, and the term $2K_q d\sigma^2/\sigma^2$ adds a contribution from gradients of uncertainty; for constant $\sigma$ the factor is a global rescaling and no region is larger than another. The fiber metric behind this example, the Fisher metric of the isotropic Gaussian family, is hyperbolic with constant curvature $-1/(2K_q)$. Hyperbolic Fisher geometry is generic for families with such symmetries~\cite{ErdmengerGrosvenorJefferson2020}, so a geometry obtained from Fisher information is not by itself evidence of physics.

The bundle-pullback formulation is the sense in which the model has the structure Wheeler asks for at the level of metric data. A metric on the base is the pullback of a fiber-respecting metric on an information-bearing bundle along an agent's section; metric structure is derived from belief data, not posited; the pulled-back tensors are observer-relative because beliefs differ; and there is no privileged geometry, only the family of agent-relative geometries induced on the same base. The construction does not derive the base, its dimension, or a Lorentzian signature.

\section{The Lorentzian Signature Problem}\label{sec:signature_resolution}

The central unsolved problem for any spacetime reading is that observed spacetime has Lorentzian signature $(-,+,+,+)$, whereas the Fisher pullbacks $G_i^{(q)}, G_i^{(p)}$ are positive semi-definite, being expectations of outer products of score vectors. When a signature appears in this framework it is a postulate, not a derivation: it lives in the Lie-algebra trace form $\operatorname{tr}(A_\mu A_\nu)$ of the connection sector, not in the Fisher-Rao metric on the belief fiber, and we make no claim that free-energy minimization selects the configuration that produces it. Earlier work treats the signature as a dynamical phase and finds a Lorentzian stationary point singled out in four dimensions~\cite{Greensite1993}; nothing comparable is attempted here, and the companion theory gives the base no Lorentzian structure~\cite[Ch.~1]{DennisGaugeVFE2026}. Two caveats govern everything below. First, in Regime~I the frames are removable at a point (Proposition~\ref{prop:regime1_removability}): by the KL invariance theorem (Theorem~\ref{thm:glk_invariance}) the functional depends on frames only through the transports $\Omega_{ij}$, and the change of reference $U_i(c) \mapsto U_i(c)g(c)$ leaves every $\Omega_{ij}$ unchanged while it can set $\operatorname{tr}(A^{(i)}_\mu A^{(i)}_\nu)$ to zero for any chosen agent (Proposition~\ref{prop:regulator_obstruction}(b)). The frame-twist form is therefore a property of a declared reference section and has no gauge-invariant content. Second, the signature results concern the frame-twist block alone, while the clock, the sectors, perceived space, and the gravity reading are readings of the positive semi-definite Fisher block.

\paragraph{Sector split.} The structure group is enlarged from $\mathrm{GL}(K_q,\mathbb{R})$ to $\mathrm{GL}(K_q,\mathbb{C})$ on the connection sector only. The Gaussian belief fiber remains a real distribution on $\mathbb{R}^{K_q}$, and belief transport $\rho(\Omega)\cdot(\mu,\Sigma) = (\Omega\mu, \Omega\Sigma\Omega^\top)$ stays restricted to $\Omega \in \mathrm{GL}(K_q,\mathbb{R})$: a genuinely complex $\Omega$ would render $\Omega\Sigma\Omega^\top$ non-Hermitian and the Gaussian relative entropy complex-valued, so real Gaussian probability, Fisher-Rao geometry, nonnegativity of relative entropy, and $\mathcal{F} \ge 0$ all live on the real sector. The complexification enters only through the frame field $\phi: \mathcal{C} \to \mathfrak{gl}(K_q,\mathbb{C})$. The construction stays inside Regime~I ($F_{\mu\nu} = 0$ throughout), so the signature mechanism runs through the pure-gauge frame-twist form and not through any curvature invariant. A stronger route that replaces the probability fiber itself by a complex or quantum statistical object (density matrices, Born-rule amplitudes, signed measures) would require a quantum-information formulation of the variational principle and is listed as future work in Section~\ref{sec:open_problems}.

\paragraph{Where an indefinite form can live.} Compact structure groups such as $\mathrm{SO}(N)$ force Riemannian geometry on the fiber, since $\Omega\Sigma\Omega^\top$ is then a similarity transform preserving eigenvalues. Non-compactness alone is necessary but not sufficient: for real invertible $\Omega$ and positive-definite $\Sigma$ the congruence $\Omega\Sigma\Omega^\top$ stays positive definite by Sylvester's law of inertia, so real $\mathrm{GL}(K_q,\mathbb{R})$ on real fibers still produces positive-definite transport. The signature resides in the trace form $\operatorname{tr}(A_\mu A_\nu)$, which on $\mathfrak{gl}(K_q,\mathbb{R})$ is already indefinite, negative definite on the skew subalgebra $\mathfrak{so}(K_q)$ and positive definite on the symmetric complement, with signature $(K_q(K_q+1)/2, K_q(K_q-1)/2)$. An indefinite restriction is therefore available on every $\mathfrak{gl}(K_q,\mathbb{R})$ with $K_q \ge 2$; what the constructions below add is an explicit frame field that produces a chosen split.

\subsection{The Postulates for an Indefinite Frame-Twist Form}

A Lorentzian-compatible signature on the horizontal block follows from the postulates below, each a postulate and not a consequence of free-energy minimization; register entry P4 summarizes the set. The indefinite form is a pullback only in the sense of the complex-bilinear trace form pulled back along the frame map; it is distinct from the positive semi-definite Fisher pullback of Eq.~\eqref{eq:induced_metric_full}. Postulates~\ref{post:pairing}, \ref{post:imaginary}, and~\ref{post:realpart} belong to the complex route only, and the all-real construction of Proposition~\ref{prop:4d_signature} dispenses with them. Two further conditions are not independent postulates: flatness is automatic once the frame is confined to an abelian subalgebra (Postulate~\ref{post:collapse}), since then $A_\mu = \partial_\mu\phi$ and $F_{\mu\nu} = 0$ identically, and a separable ansatz for the frame components only diagonalizes the chart, as the worked example shows.

\renewcommand{\thepostulate}{P4.\arabic{postulate}}
\begin{postulate}[Generator with nonzero trace-square]\label{post:noncompact}
The temporal and spatial frame directions use generators $T$ with $\operatorname{tr}(T^2) \neq 0$; in the two-dimensional example, the symmetric traceless $T = \operatorname{diag}(1,-1) \in \mathfrak{sl}(K_q,\mathbb{R})$ with $\operatorname{tr}(T^2) = 2$. A compact generator ($\operatorname{tr}(T^2) < 0$) exchanges which direction is timelike, and only a nilpotent generator ($\operatorname{tr}(T^2) = 0$) removes the form.
\end{postulate}
\begin{postulate}[Complex-bilinear pairing; complex route only]\label{post:pairing}
The form $\langle A, B\rangle = \operatorname{tr}(AB)$ on $\mathfrak{gl}(K_q,\mathbb{C})$, in place of the sesquilinear Frobenius pairing $\operatorname{tr}(AB^\dagger)$, which is positive definite and would give no sign change; the bilinear, non-Hermitian character is load-bearing.
\end{postulate}
\begin{postulate}[Imaginary temporal frame; complex route only]\label{post:imaginary}
The assignment $\phi_0 = i\psi_0 T$ along the distinguished base direction. This is an imaginary coefficient in the Lie algebra; it is not a Wick rotation, which continues a base coordinate and never meets a real-part projection.
\end{postulate}
\begin{postulate}[Real-part projection; complex route only]\label{post:realpart}
Extraction of $G^{\mathrm{Lor}}_{\mu\nu} := \operatorname{Re}(G_{\mu\nu})$, a rank-changing operation (rank one complex to rank two real) for which no principle in the construction mandates discarding the imaginary off-diagonal part.
\end{postulate}
\begin{postulate}[Single-generator collapse]\label{post:collapse}
$T_0 = T_1 = T$, confining $\phi$ to the abelian subalgebra $\mathbb{C}T$.
\end{postulate}
\begin{postulate}[Mutual trace-orthogonality of the spatial generators]\label{post:traceorth}
In any multi-generator extension the spatial generators are mutually trace-orthogonal, $\operatorname{tr}(T_a T_b) = 0$ for $a \neq b$, so that the spatial block of the frame-twist form is diagonal at the abelian-exact order; the commuting disjoint-support generators of Proposition~\ref{prop:4d_signature} satisfy this by construction.
\end{postulate}

Each step is mathematically well defined. In the two-dimensional example the complex route is forced by abelian exactness: a single real generator has a trace form of one sign, so within a one-generator abelian subalgebra a sign change can enter only through $i^2 = -1$. The indefiniteness then comes from Postulates~\ref{post:pairing}, \ref{post:imaginary}, and~\ref{post:realpart} together, since $\operatorname{Re}((ia+b)^2) = b^2 - a^2$ is the hyperbolic quadratic form on $\mathbb{C} = \mathbb{R}^2$, and choosing which base direction receives the factor $i$ is choosing which direction is timelike. A real frame mixing skew and symmetric generators reaches an indefinite form without complexification, and Proposition~\ref{prop:4d_signature} carries this out in four dimensions.

\begin{proposition}[Conditional $(1{+}1)$ signature]\label{prop:conditional_signature}
Under Postulates~\ref{post:noncompact} to~\ref{post:collapse} on a two-dimensional base, the real part of the frame-twist form $\operatorname{tr}(A_\mu A_\nu)$ has rank two and signature $(-,+)$ wherever $d\psi_0$ and $d\psi_1$ are linearly independent, and at each such point the orthonormal-frame transformations form $\mathrm{O}(1,1)$, whose identity component is $\mathrm{SO}^+(1,1)$.
\end{proposition}
\begin{proof}
The computation of Section~\ref{sec:worked_signature}.
\end{proof}
The proposition is a conditional result with status P under its postulates; any physical reading of its conclusion is status S, and whether the dynamics select its hypotheses is open.

\subsection{Worked Example: \texorpdfstring{$\mathrm{GL}(2,\mathbb{C})$}{GL(2,C)} Frames and \texorpdfstring{$(1{+}1)$}{(1+1)} Lorentzian Signature}
\label{sec:worked_signature}

On a two-dimensional base $\mathcal{C}$ with coordinates $(x^0, x^1)$, take $T = \operatorname{diag}(1,-1)$ with $\operatorname{tr}(T^2) = 2$ (Postulate~\ref{post:noncompact}) and the imaginary temporal component of Postulate~\ref{post:imaginary},
\begin{equation}
\phi(x^0, x^1) = i\psi_0(x^0, x^1) T + \psi_1(x^0, x^1) T.
\label{eq:complex_gauge_frame}
\end{equation}
Under the single-generator collapse (Postulate~\ref{post:collapse}), $\phi$ lies in the one-dimensional abelian subalgebra $\mathbb{C}T$, so the connection is $A_\mu = \partial_\mu\phi$ exactly (the example is abelian-exact, not linearized) and $F_{\mu\nu} = 0$. For general components, $G_{\mu\nu} = \operatorname{tr}(A_\mu A_\nu) = 2(i\partial_\mu\psi_0 + \partial_\mu\psi_1)(i\partial_\nu\psi_0 + \partial_\nu\psi_1)$, a complex form of rank one with vanishing determinant, and
\begin{equation}
\operatorname{Re}G = 2\left(d\psi_1 \otimes d\psi_1 - d\psi_0 \otimes d\psi_0\right), \qquad \det \operatorname{Re}G = -4\mathcal{J}^2,
\end{equation}
with $\mathcal{J}$ the Jacobian determinant of $(\psi_0, \psi_1)$, so the signature is $(-,+)$ wherever $\mathcal{J} \neq 0$. In the separable chart $\psi_0 = \psi_0(x^0)$, $\psi_1 = \psi_1(x^1)$ the components are
\begin{align}
G_{00} &= i^2 (\partial_0 \psi_0)^2 \operatorname{tr}(T^2) = -2(\partial_0 \psi_0)^2 < 0, \\
G_{11} &= (\partial_1 \psi_1)^2 \operatorname{tr}(T^2) = +2(\partial_1 \psi_1)^2 > 0, \\
G_{01} &= i(\partial_0 \psi_0)(\partial_1 \psi_1)\operatorname{tr}(T^2) \in i\mathbb{R},
\end{align}
and the real-part projection $G^{\mathrm{Lor}}_{\mu\nu} := \operatorname{Re}(G_{\mu\nu})$ sets $G^{\mathrm{Lor}}_{01} = 0$ and yields the rank-two real metric
\begin{equation}
\boxed{ds^2 = -2(\partial_0 \psi_0)^2 (dx^0)^2 + 2(\partial_1 \psi_1)^2 (dx^1)^2 .}
\label{eq:lorentzian_metric}
\end{equation}
Separability therefore only diagonalizes the chart: a non-separable frame gives, for example, $\operatorname{Re}G_{00} = +6$ with eigenvalues $-0.081$ and $7.901$, still Lorentzian. With a compact generator $T = J$ ($\operatorname{tr}(J^2) = -2$) the real part is $2(d\psi_0 \otimes d\psi_0 - d\psi_1 \otimes d\psi_1)$, again of signature $(1,1)$ with the timelike role moved to $x^1$, and with a nilpotent generator the form vanishes. Calling Eq.~\eqref{eq:lorentzian_metric} a physical metric requires both the imaginary frame of Postulate~\ref{post:imaginary} and the real-part projection of Postulate~\ref{post:realpart}, and this is a derivation gap: writing $\chi = i\psi_0 + \psi_1$, the projected form is the difference of squares of the differentials of two scalar functions, obtainable on any manifold with no gauge theory, so the gauge apparatus contributes only the normalization $\operatorname{tr}(T^2)$ and the sign conventions. The $\pm 2$ coefficients are artifacts of the unnormalized $T$; under $\operatorname{tr}(T^2) = 1$ they become $\pm 1$, and the construction fixes the metric only up to a conformal class. At each non-degenerate point the orthonormal-frame transformations form $\mathrm{O}(1,1)$, whose identity component $\mathrm{SO}^+(1,1)$ is the group of $(1{+}1)$ boosts, as the relation $\Lambda^\top\eta\Lambda = \eta$ with $\eta = \operatorname{diag}(-1,+1)$ shows.

\begin{proposition}[Rank obstruction of the single-generator collapse]\label{prop:rank_obstruction}
The single-generator collapse does not extend to four dimensions: with one shared generator and $\operatorname{tr}(T^2) = 2$ the frame-twist form is a rank-one outer product whose real-part projection raises the rank only to two, with eigenvalues $\{-2(\partial_0\psi_0)^2, 2[(\partial_1\psi_1)^2 + (\partial_2\psi_2)^2 + (\partial_3\psi_3)^2], 0, 0\}$ in the separable chart and degenerate signature $(-,+,0,0)$, whose non-degenerate part again carries only $\mathrm{O}(1,1)$.
\end{proposition}

A rank-four non-degenerate $(-,+,+,+)$ form, with local frame group $\mathrm{O}(1,3)$ and identity component $\mathrm{SO}^+(1,3)$, requires abandoning the collapse in favor of four independent generators, and complexification is dispensable for this existence question: the economical choice is real.

\begin{proposition}[Four-dimensional all-real signature construction]\label{prop:4d_signature}
Let $K_q \geq 6$. Choose a real skew generator $S$ supported on two coordinates with $\operatorname{tr}(S^2) = -2$, and three traceless, mutually trace-orthogonal real diagonal generators $T_1, T_2, T_3$ supported on the remaining coordinates (Postulate~\ref{post:traceorth}). All four generators commute, so $\phi = \psi_0 S + \sum_a \psi_a T_a$ lies in an abelian subalgebra, $A_\mu = \partial_\mu\phi$ exactly (Regime~I, $F_{\mu\nu} = 0$), and the construction stays inside $\mathrm{GL}(K_q,\mathbb{R})$. Because $\operatorname{tr}(ST) = 0$ identically for skew $S$ and symmetric $T$, all temporal-spatial cross terms vanish without projection, and the frame-twist form $G_{\mu\nu} = \operatorname{tr}(A_\mu A_\nu)$ is real, of rank four and signature $(-,+,+,+)$ wherever the four differentials $d\psi_0, d\psi_1, d\psi_2, d\psi_3$ are linearly independent, with orthonormal-frame group $\mathrm{O}(1,3)$ at each such point. No imaginary unit, complex pairing, or real-part projection enters: Postulates~\ref{post:pairing}, \ref{post:imaginary}, and~\ref{post:realpart} are not required.
\end{proposition}
\begin{proof}
$S$ has support disjoint from the $T_a$ and so commutes with them, and the $T_a$ commute with one another because they are diagonal; hence $\phi$ is abelian and $A_\mu = \partial_\mu\phi$ exactly. Normalizing $\operatorname{tr}(T_a T_b) = 2\delta_{ab}$ gives $G_{\mu\nu} = -2(\partial_\mu\psi_0)(\partial_\nu\psi_0) + 2\sum_a (\partial_\mu\psi_a)(\partial_\nu\psi_a)$. For skew $S$ and symmetric $T$, $\operatorname{tr}(ST) = \operatorname{tr}((ST)^\top) = -\operatorname{tr}(ST)$ forces $\operatorname{tr}(ST) = 0$, so no cross terms survive. Thus $G = 2\psi^*\eta$ with $\eta = \operatorname{diag}(-1,1,1,1)$, and where $d\psi$ is invertible Sylvester's law of inertia gives the stated signature.
\end{proof}

Several limits bound what the proposition delivers. The metric is $G = 2\psi^*\eta$, the pullback of the Minkowski metric along $\psi = (\psi_0, \psi_1, \psi_2, \psi_3): \mathcal{C} \to \mathbb{R}^{1,3}$, hence coordinate-flat: the proposition exhibits signature availability, not curvature, and the derivation-gap critique of the two-dimensional example transfers, since the gauge apparatus again contributes only the trace normalization. The signature is fixed by the choice of one skew and three diagonal generators, which is the content of P4, so availability is immediate from the indefinite trace form and selection of the $1{+}3$ split (C4 of Section~\ref{sec:open_problems}) is the open problem. The bound $K_q \geq 6$ is the minimum under the disjoint-support hypothesis with traceless spatial generators, and $K_q \geq 5$ suffices under that hypothesis without tracelessness; without disjoint supports, commuting trace-orthogonal sets exist at $K_q = 5$ with traceless spatial generators (for instance $S = J \oplus 0$ with $T_a = a_a I_2 \oplus \operatorname{diag}(c_a, d_a, e_a)$, $2a_a + c_a + d_a + e_a = 0$) and at $K_q = 4$ without tracelessness. The irreducible content is the generator split: one compact and three non-compact directions, mutually trace-orthogonal and commuting.

As stated at the head of this section, the signature is a property of a declared reference section. A gauge-invariant candidate exists in the relational form $\operatorname{tr}((A_i - A_j)_\mu(A_i - A_j)_\nu) = \operatorname{tr}((\Omega_{ij}^{-1}\partial_\mu\Omega_{ij})(\Omega_{ij}^{-1}\partial_\nu\Omega_{ij}))$ on agent pairs, which is indefinite on $\mathfrak{gl}(K_q,\mathbb{R})$ for the same reason as $\operatorname{tr}(A_\mu A_\nu)$; it requires base derivatives of frames and is deferred to Supplement~B. Because the indefinite form lives on the horizontal block alone while the clock, the sectors, and the gravity reading of Section~\ref{sec:gravity_pullback} are readings of the Fisher block, the signature that matters for them is that of the combined tensor of Eq.~\eqref{eq:combined_tensor}, which depends on the relative size of the two blocks: in the four-generator frame with one Gaussian mean $\mu = a\psi_0$ of variance $\sigma^2$, the combined tensor in the chart $\psi$ is $\operatorname{diag}(-2\kappa + a^2/\sigma^2, 2\kappa, 2\kappa, 2\kappa)$, Riemannian whenever $a^2/\sigma^2 > 2\kappa$. We do not identify a single tensor that is both Lorentzian and the carrier of the clock, the sectors, and the gravity reading. The open problem accordingly sharpens: not whether an indefinite frame-twist form can be exhibited in four dimensions, but whether a gauge-invariant form can be, and whether free-energy dynamics select one compact and three non-compact commuting frame directions.

\subsection{Alternative Route: Lorentzian Conformal Class from Finite-Speed Epistemic Causality}
\label{sec:causal_cone_route}

A second route reaches a Lorentzian form without complexifying the connection, using a disjoint postulate set (register entry P5). It assumes a coarse-grained shared update parameter $t$ used as a base coordinate, a positive-definite spatial information metric $h_{ab}$, and a finite maximal propagation speed $c_\mathcal{I}$ for epistemic perturbations, with the region of possible influence the $h$-ball $h_{ab}dx^a dx^b \le c_\mathcal{I}^2 dt^2$. The boundary of possible influence is then the null cone of $g_{\mu\nu} dc^\mu dc^\nu = -c_\mathcal{I}^2 dt^2 + h_{ab} dx^a dx^b$, which is of Lorentzian signature $(-,+,\ldots,+)$ by inspection and of signature $(-,+,+,+)$ for three spatial dimensions. The quadratic cone requires the isotropy assumption; a direction-dependent speed bound gives a Finsler cone. The shared parameter $t$ is a preferred foliation, so Lorentz boosts are symmetries of the cone but not of the first-order dynamics, as in analog gravity~\cite{BarceloLiberatiVisser2011}. The cone fixes $g$ only up to a conformal factor and does not select the dimension count. This is as it must be: the causal-order theorems cited in Section~\ref{sec:pullback}~\cite{HawkingKingMcCarthy1976, Malament1977} fix the metric only up to a conformal factor and presuppose the signature, so they do not produce one from a non-Lorentzian order, and causal set theory supplies the missing scale by counting elements~\cite{BombelliLeeMeyerSorkin1987}. The speed $c_\mathcal{I}$ is postulated and carries no units. The finite-speed postulate is in tension with the first-order natural-gradient dynamics only in the continuum limit, where the naive limit is parabolic with infinite signal speed; on a bounded-degree agent graph, first-order dynamics already has an effective light cone of Lieb-Robinson type: for linear dynamics $\dot{x} = -Lx$ with a unit-weight graph Laplacian of maximal degree $\Delta$, $(L^k)_{xy}$ vanishes for $k$ below the graph distance $d(x,y)$, so $|(e^{-tL})_{xy}| \le \sum_{k \ge d(x,y)}(2\Delta t)^k/k!$~\cite{LiebRobinson1972, Marchioro1978}. Relating that graph cone to a causal structure on $\mathcal{C}$ is base-transport work deferred to Supplement~B. The two routes are parallel existence statements with disjoint open problems (dynamical selection of one compact and three non-compact commuting frame directions, and gauge-invariant content, for the algebraic route; the conformal scale, the dimension count, the isotropy assumption, and the continuum limit of the dynamics for the causal route), and neither derives the signature from variational dynamics.

\subsection{Temporal Direction from Belief Trajectories}

An agent's belief trajectory singles out a line in its tangent space through the Gram-Schmidt decomposition
\begin{equation}
T_q\mathcal{B} = \mathbb{R}\dot{q} \oplus \dot{q}^\perp,
\label{eq:tangent_decomp}
\end{equation}
but assigns that line neither a negative norm nor an orientation; the orientation is inherited from the flow parameter $t$, equivalently from the decrease of $\mathcal{F}$ along the flow. Writing
\begin{equation}
ds^2 = -g_{\mathcal{B}}(\dot{q}, \dot{q}) dt^2 + \sum_i g_{\mathcal{B}}(e_i, e_i) dx_i^2 = -d\vartheta^2 + \sum_i g_{\mathcal{B}}(e_i, e_i) dx_i^2,
\label{eq:lor_belief_metric}
\end{equation}
the leading minus sign is imposed by ansatz on the fiber tangent space $T_q\mathcal{B}$, and no construction in this paper produces it. The signature constructions of Section~\ref{sec:signature_resolution} give indefinite forms on the base tangent space built from the connection, not a sign on $g_{\mathcal{B}}$, and the relative entropy between real Gaussian beliefs is unaffected by any real-versus-imaginary choice in the present formulation. Equation~\eqref{eq:lor_belief_metric} records what a Lorentzian reading of the belief trajectory would need; it is not a derivation.

The boxed item below records what is derived and what is postulated.

\begin{center}
\fbox{\begin{minipage}{0.9\columnwidth}
\textbf{Belief Tangent-Space Decomposition and the Spacetime Reading: What Is Derived and What Is Postulated}

\textbf{What is derived:}
A belief trajectory $q(t)$ through the fiber $\mathcal{B}$ has tangent $\dot{q} \in T_q\mathcal{B}$ with the Gram-Schmidt decomposition
\[
T_q\mathcal{B} = \underbrace{\mathbb{R}\dot{q}}_{\substack{\text{trajectory} \\ \text{1 dim}}}
\oplus \underbrace{\dot{q}^\perp}_{\substack{\text{orthogonal} \\ \text{complement}}}
\]
of Eq.~\eqref{eq:tangent_decomp}, which singles out the trajectory line without assigning it a sign or an orientation, together with the spectral decomposition of the pullback tensor $G_i(c)$ on the base tangent space $T_c\mathcal{C}$ relative to a declared background chart metric (Section~\ref{sec:observable_sectors}). Time and the sectors live on the base: under the identification $q(t) = \sigma_i(c(t))$ of the trajectory with the section evaluated along a base curve, the temporal direction is the tangent $\dot{c}$ of that curve, whose pushforward $d\sigma_i(\dot{c}) = \dot{q}$ recovers the fiber trajectory direction, and the observable, subthreshold, and internal sectors $\mathcal{D}_{\text{obs}}, \mathcal{D}_{\text{subthresh}}, \mathcal{D}_{\text{internal}}$ are eigen-direction classes of $G_i(c)$ on $T_c\mathcal{C}$, not of any operator on $\dot{q}^\perp$, on which no eigenvalue problem is defined.

\textbf{What is postulated:} (i) the eigenvalue hierarchy $\lambda_{\text{obs}} \gg \lambda_{\text{subthresh}} \gg \lambda_{\text{internal}} \approx 0$ on the spectrum of $G_i(c)$, including the choice of four large eigenvalues in $\mathcal{D}_{\text{obs}}$, one temporal and three spatial (register entry P6), together with any three-way refinement of $\dot{q}^\perp$ into fiber subspaces $V_{\text{obs}} \oplus V_{\text{subthresh}} \oplus V_{\text{internal}}$ mirroring the base sectors, with $\mathbb{R}\dot{q}$ mirroring the temporal and $V_{\text{obs}}$ the three spatial directions of $\mathcal{D}_{\text{obs}}$; (ii) Lorentzian signature $(-,+,+,+)$ on $\mathbb{R}\dot{q} \oplus V_{\text{obs}}$, as the ansatz of Eq.~\eqref{eq:lor_belief_metric} with no construction behind it; (iii) on the complex route, Postulates~\ref{post:pairing}, \ref{post:imaginary}, and~\ref{post:realpart} of Section~\ref{sec:signature_resolution} (the all-real construction of Proposition~\ref{prop:4d_signature} dispenses with them), acting on the frame-twist block of a declared reference section; (iv) in any multi-generator extension, the mutual trace-orthogonality of the spatial generators, Postulate~\ref{post:traceorth}, exact at the abelian-exact order and satisfied by commuting disjoint-support generators in Proposition~\ref{prop:4d_signature}, while beyond that order the $\mathrm{dexp}$ expansion of $A_a = U^{-1}\partial_a U$ mixes generators and exact diagonality would additionally require commuting frame components; (v) that one tensor carries both the signature and the clock and sectors, which no construction here shows. None of (i) to (v) is derived from free-energy minimization in the present formulation.

\textbf{An interpretive reading:} The above structure permits a three-sector reading consisting of a $(1+3)$ spacetime sector (temporal plus observable spatial), a subthreshold informational sector with moderate squared information rate, and an internal sector where that rate is negligible. We do not claim to have derived physical spacetime, the Lorentz signature, or the dimension count from the framework; the construction shows that the framework is compatible with such a reading under the postulates listed above.
\end{minipage}}
\end{center}
